\chapter{Getting Started}%
\label{ch:getting-started} % chktex 8

In this chapter we take the very first steps: we get the Salam compiler ready,
write our first program, and run it. Do not worry if you have never written a
line of code before this chapter takes you by the hand and walks you forward
one step at a time.

\section{What is a program?}

A computer is a machine that is extraordinarily fast and extraordinarily
literal. It does not know what you want; you must \emph{tell} it, in exact
detail. A list of instructions, written in order so that the computer can carry
out a task, is called a \textbf{program}. The language we write those
instructions in is called a \textbf{programming language}.

Salam is one such language, with one special quality: you can write your
instructions in \emph{Persian} as comfortably as in English. The instruction
``print something on the screen'' is written \code{println "Salam"} in English,
or \code{\fa{چاپ} "\fa{سلام}"} in Persian. Both mean exactly the same thing to
the computer.

\begin{note}
A computer does \emph{precisely} what you tell it no more, no less. So when
a program does not behave the way you expected, the mistake is almost always in
the instructions you gave, not in the machine. This is completely normal; even
professional programmers run into it every single day. Learning to read what the
computer actually did, rather than what you \emph{meant}, is half the craft.
\end{note}

\section{What is a compiler, and why do we need one?}

The program you write is meant to be readable by a \emph{human}. The computer,
however, does not understand it directly deep down it only understands tiny
numeric instructions. We need a program that translates what we wrote into
something the machine can run. That translator is called a \textbf{compiler}.
The compiler for the Salam language is called \code{salam}.

\code{salam} does a great deal of work for us: it reads the text of your
program, checks that it makes sense, reports any mistakes, and finally produces
something you can run.

\begin{deep}[In Depth: the journey of a program]
When \code{salam} processes your program, it passes through several stages:
\emph{lexing} (breaking the text into tokens), \emph{parsing} (building a tree
that represents the structure), \emph{semantic analysis} (checking types and
names), and finally \emph{code generation}. Salam's main backend translates your
program into C source code and then invokes a C compiler (\code{tcc} by
default, but \code{gcc} or \code{clang} work too) to produce a native
executable. There is also a built-in \emph{interpreter} that can run your % chktex 8
program directly without any C toolchain, and an experimental LLVM backend. We
will refer to these stages occasionally, but you never need to know their
details to write programs.
\end{deep}

\section{Getting the compiler}

To work with Salam you need the \code{salam} compiler. It is built from source.
In the root folder of the Salam project, run:

\begin{shell}
sh tools/build-compiler.sh
\end{shell}

When it finishes, an executable called \code{salam} (or \code{salam.exe} on
Windows) appears in the project folder. Alternatively you can build with CMake,
which also builds the runtime library and enables the test suite:

\begin{shell}
cmake -B build && cmake --build build
\end{shell}

\begin{tip}
To confirm the compiler is working, ask it for its version:
\begin{shell}
$ salam version
salam 0.2.6
\end{shell}
\end{tip}

\begin{note}
Throughout this book we assume \code{salam} is reachable from your terminal:
either it sits in your current folder (call it as \code{./salam}) or it is
installed on your system \code{PATH} (call it as \code{salam}). For brevity we
will simply write \code{salam} everywhere.
\end{note}

\section{Your first program: Hello, Salam!}

There is a long-standing tradition among programmers: the very first program you % chktex 8
write in a new language is one that displays the words ``Hello, World!'' Let us
honour that tradition.

Create a file named \code{hello.salam} and type these few lines into it:

\begin{salam}
func main:
    name : str = "Salam"
    println "Hello,", name
    println "2 + 3 =", 2 + 3
end
\end{salam}

\begin{note}
Every Salam program file ends in \code{.salam}. The extension tells the tools
(and you) that the file contains Salam source code.
\end{note}

\subsection{Running the program}

The simplest way to see the result is to use Salam's built-in interpreter, which % chktex 8
runs your program directly no C compiler required. In your terminal, type:

\begin{shell}
salam exec hello.salam
\end{shell}

and the output will be:

\begin{progout}
Hello, Salam
2 + 3 = 5
\end{progout}

If you would rather build a standalone executable (which runs faster and is what
you ship to other people), use the \code{build} command:

\begin{shell}
$ salam build hello.salam --output=hello.exe
$ ./hello.exe
Hello, Salam
2 + 3 = 5
\end{shell}

There is also a one-shot \code{run} command that builds and immediately runs % chktex 8
without leaving any files behind:

\begin{shell}
salam run hello.salam
\end{shell}

\begin{tip}
What is the difference between \code{exec}, \code{run}, and \code{build}?
\code{exec} runs the program \emph{immediately} with the interpreter and needs
no C toolchain it is perfect for learning and quick experiments.
\code{build} produces a native executable you can run again and again, and is
what you use to deliver a finished program. \code{run} is the convenient middle % chktex 12
ground: compile and run once, keep nothing. Throughout most of this book we will
use \code{exec} for simplicity.
\end{tip}

\section{Anatomy of the first program}

Let us now go through the program line by line. Do not skim any part of this ---
these few lines are the doorway to every concept that follows.

\paragraph{Line 1 the start of the program.}
\begin{salam}
func main:
\end{salam}
This line says: ``define a \textbf{function} named \textbf{main}.'' A function is
a named bundle of instructions. The function named \code{main} is the
\emph{starting point} of every Salam program: when you run the program, the
computer begins here. The colon (\code{:}) at the end of the line says that the
body of the function starts on the next line.

\paragraph{Line 2 a variable.}
\begin{salam}
    name : str = "Salam"
\end{salam}
Here we create a \textbf{variable} named \code{name}. Picture a variable as a
labelled box that holds a value. In this box we have placed the text
\code{"Salam"}. The word \code{str} states that the box holds \emph{text} (more
on types in the next chapter). Any piece of text is written between double
quotes (\code{"\,}). Notice that this line begins with a few spaces of
\textbf{indentation}; that whitespace shows the line lives \emph{inside} the
\code{main} function.

\paragraph{Line 3 printing output.}
\begin{salam}
    println "Hello,", name
\end{salam}
The word \code{println} (``print line'') tells the computer to display something
on the screen and then move to the next line. Here we display two things: the
fixed text \code{"Hello,"} and then the contents of the variable \code{name}
(which is \code{"Salam"}). The two are separated by a comma (\code{,}), and Salam
puts a single space between them. That is why the output reads
\code{Hello, Salam}.

\paragraph{Line 4 a little arithmetic.}
\begin{salam}
    println "2 + 3 =", 2 + 3
\end{salam}
Salam does not only print text; it can calculate too. The expression
\code{2 + 3} is worked out \emph{before} printing, and the result (\code{5}) is
displayed. Notice the difference: \code{"2 + 3 ="} is inside quotes, so it is a
piece of \emph{text} printed exactly as written; but \code{2 + 3} is outside
quotes, so it is a \emph{calculation}.

\begin{warn}
Text in quotes is printed verbatim; code outside quotes is evaluated. Writing
\code{println "2 + 3"} would print the three characters \texttt{2 + 3}, not the
number \texttt{5}. The quotes make all the difference.
\end{warn}

\paragraph{Last line the end of the function.}
\begin{salam}
end
\end{salam}
The word \code{end} tells Salam that the body of \code{main} is finished. Any
block that opens with a colon (\code{:}) must be closed with \code{end}. This is
a rule you will see throughout the language: \emph{blocks open with a colon and
close with} \code{end}.

\section{The same program in Persian}

One of the joys of Salam is that any program can be written in Persian too. This
is the exact Persian equivalent of the program above:

\begin{salamfa}
تابع اصلی:
    نام : رشته = "سلام"
    چاپ "سلام،", نام
    چاپ "2 + 3 =", 2 + 3
تمام
\end{salamfa}

The structure is identical; only the words changed:
\code{func}~$\leftrightarrow$~\code{\fa{تابع}},
\code{main}~$\leftrightarrow$~\code{\fa{اصلی}},
\code{str}~$\leftrightarrow$~\code{\fa{رشته}},
\code{println}~$\leftrightarrow$~\code{\fa{چاپ}}, and
\code{end}~$\leftrightarrow$~\code{\fa{تمام}}. A full table of these
equivalents is in Appendix~A. To compile a Persian file, tell the compiler the
language:

\begin{shell}
salam exec hello_fa.salam --lang=fa
\end{shell}

\begin{note}
The entry function's name is itself part of the language: in English it is
\code{main}, in Persian \code{\fa{اصلی}}. We will return to bilingual
programming properly in Chapter~22. For now, just notice that the \emph{ideas}
are the same in both languages only the spelling of the keywords differs.
\end{note}

\section{Comments: notes to yourself}

A \textbf{comment} is text that the compiler ignores completely. Comments are
written for \emph{humans} to explain what a piece of code does, or to leave a
reminder. A line comment starts with \code{//} and runs to the end of the line.
A block comment is wrapped in \code{/* ... */} and may span several lines.

\begin{salam}
func main:
    // This is a line comment: the compiler skips it.
    println "working"          // comments can also follow code

    /* This is a block comment.
       It can stretch across
       as many lines as you like. */
    println "done"
end
\end{salam}

\begin{progout}
working
done
\end{progout}

\begin{tip}
Good comments explain \emph{why}, not \emph{what}. The code already says what it
does; a comment is most valuable when it records the reasoning that the code
cannot show by itself.
\end{tip}

\section{Printing with and without a new line}

Salam has two main ways to display output:

\begin{itemize}[leftmargin=1.4em]
  \item \code{println} displays its arguments and \emph{then moves to the
        next line}.
  \item \code{print} displays its arguments but \emph{stays on the same
        line}; whatever is printed next continues right after it.
\end{itemize}

Look at this example:

\begin{salam}
func main:
    print "Salam "
    print "world"
    println ""
    println "a fresh line"
end
\end{salam}

\begin{progout}
Salam world
a fresh line
\end{progout}

The two \code{print} calls both wrote on the same line, and then
\code{println ""} (which prints an empty piece of text) moved us down to the
next line. This simple trick writing with \code{print} and ending the line
with \code{println ""} will be used again and again in the chapter on loops
to draw shapes.

\begin{deep}[In Depth: writing the arguments of print]
\code{print} and \code{println} are statements, not ordinary functions, so
their arguments are written \textbf{without} wrapping parentheses and
separated by commas: \code{println "hi"} and \code{println name, age}.
Wrapping the whole argument list in parentheses is an error, and Salam will
refuse to compile it. Parentheses inside an argument are of course still
fine, since there they group an expression rather than wrap the list:
\code{println 1, (3 * 4)} and \code{println (2 + 3) * 4} are both correct.
Blocks, likewise, always open with a colon and close with \code{end}; there % chktex 8
is no curly-brace alternative. Remember the \code{salam format} command,
which reformats your code into the canonical style automatically.
\end{deep}

\section{Starting a new project}

When you begin something larger than a single file, the compiler can scaffold a
project folder for you:

\begin{shell}
salam new myproject --lang=en
\end{shell}

This creates a \code{myproject} directory with a starter program already in
place, so you can build and run it immediately. For the small examples in the
coming chapters, a single \code{.salam} file is all you need.

\section{Chapter summary}

In this chapter we learned that:
\begin{itemize}[leftmargin=1.4em]
  \item A program is a list of instructions, and the compiler \code{salam}
        translates it for the computer.
  \item Every Salam program starts at the \code{main} function
        (\code{\fa{اصلی}} in Persian).
  \item Blocks open with a colon (\code{:}) and close with \code{end}.
  \item \code{println} prints and moves to a new line; \code{print} prints and
        stays on the line.
  \item Text goes in double quotes; code outside quotes is evaluated.
  \item Comments (\code{//} and \code{/* ... */}) are ignored by the compiler.
  \item \code{salam exec} runs a program instantly; \code{salam build} makes a
        standalone executable.
  \item Any program can be written in Persian or English with identical meaning.
\end{itemize}

\section{Exercises}

\exercise Write a program that prints your name and your age on two separate
lines.

\exercise Write a program that prints the result of \code{12 * 8} together with a
short label, for example: \texttt{12 * 8 = 96}.

\exercise Write ``Hello, World!'' once in English and once in Persian, and run
both. Is the output the same?

\exercise Using \code{print} (not \code{println}), display your first and last
name on a \emph{single} line with one space between them, then move to a new
line.

\exercise Add a block comment at the top of one of your programs describing what
it does, then run it and confirm the comment changes nothing.
